\documentclass[11pt]{article}
\usepackage[a4paper,margin=28mm]{geometry}
\usepackage[T1]{fontenc}
\usepackage{amsmath,amsthm,mathtools}
\usepackage{newtxtext,newtxmath}
\usepackage{booktabs,array}
\newcolumntype{P}[1]{>{\raggedright\arraybackslash}p{#1}}
\usepackage{microtype}
\usepackage{xcolor}
\usepackage{hyperref}
\usepackage{enumitem}
\usepackage[nameinlink,noabbrev]{cleveref}

\hypersetup{colorlinks=true,linkcolor=blue!55!black,citecolor=blue!55!black,urlcolor=blue!55!black}
\setlist{nosep}
\setlength{\parindent}{1.2em}
\setlength{\parskip}{0pt}
\setlength{\textfloatsep}{14pt plus 2pt minus 3pt}
\setlength{\floatsep}{11pt plus 2pt minus 2pt}
\emergencystretch=2em

\newtheorem{theorem}{Theorem}[section]
\newtheorem{lemma}[theorem]{Lemma}
\newtheorem{proposition}[theorem]{Proposition}
\newtheorem{corollary}[theorem]{Corollary}
\theoremstyle{definition}
\newtheorem{definition}[theorem]{Definition}
\newtheorem{remark}[theorem]{Remark}
\newtheorem{fact}[theorem]{Conjecture}
\crefname{fact}{Conjecture}{Conjectures}

\newcommand{\rad}{\operatorname{rad}}
\newcommand{\ecc}{\operatorname{ecc}}
\newcommand{\dist}{\operatorname{dist}}
\newcommand{\diam}{\operatorname{diam}}
\newcommand{\mindeg}{\delta}
\newcommand{\NC}{\mathrm{ncv}}

\title{A Self-Contained Proof of Graffiti.pc Conjecture 84:\\
$t(G)\ge 2r(G)/\delta(G)$}

\author{Hai Luo\\[6pt]
  \normalsize Tsinghua Shenzhen International Graduate School, Tsinghua University\\
  \normalsize Shenzhen, China\\[2pt]
  \normalsize\texttt{h-luo26@mails.tsinghua.edu.cn}}

\date{September 21, 2026}

\begin{document}
\maketitle

\begin{abstract}
For a finite connected graph $G$, write $r(G)$ for the radius, $\delta(G)$ for
the minimum degree, and $t(G)$ for the largest order of an induced subtree.
Graffiti.pc Conjecture 84 (West's catalogue, 2004) asserts
\[
        t(G)\ge \frac{2r(G)}{\delta(G)}.
\]
Erd\H{o}s, Saks, and S\'os proved that every connected graph of radius $r$
contains an induced path on at least $2r-1$ vertices, which settles the
conjecture whenever $\delta(G)\ge2$.  The remaining case is the content of
this paper: the \emph{leaf lemma}, stating that every finite connected graph
with a leaf contains an induced tree on at least $2r(G)$ vertices, thereby
strengthening the induced-path bound to exactly $2r$ precisely when
$\delta(G)=1$.  We prove the leaf lemma and hence the conjecture.  The
presentation is self-contained.  In particular the Erd\H{o}s--Saks--S\'os
theorem is not used (a geodesic bound suffices for the case $\delta\ge2$),
and the radius-drop criterion together with the structure theorem for
vertex-radius-decreasing graphs $G=S_\ell(B)$---classically due to Gliviak
and Fajtlowicz---are re-proved here from elementary principles; a rooted
form of Chung's induced-path construction is written out in full, including
two details that are usually left implicit.  An accompanying Lean~4/Mathlib
development verifies the leaf lemma and the final theorem under a pinned
toolchain; the formalization is a supplementary verification layer.
\end{abstract}

\medskip
\noindent\textbf{Keywords.} induced trees; graph radius; Graffiti.pc
conjectures; cut vertices; radius-critical graphs; Lean 4.

\section{Introduction}\label{sec:intro}

All graphs are finite, simple, and undirected.  For a connected graph $G$
write
\[
 \ecc_G(v)=\max_{x\in V(G)} d_G(v,x),\qquad r(G)=\min_{v\in V(G)}\ecc_G(v),
\]
and let $\delta(G)$ be the minimum degree.  Let $t(G)$ be the maximum
cardinality of a vertex set inducing a tree.  A \emph{leaf} is a vertex of
degree one, $p(u)$ denotes the neighbour of a leaf $u$, and $v$ is a
\emph{cut vertex} if $G-v$ is disconnected; otherwise $v$ is a non-cut vertex.

\begin{fact}[Graffiti.pc \#84; West~\cite{WestGraffiti}]\label{fact:84}
If $G$ is connected, then $t(G)\ge 2r(G)/\delta(G)$.
\end{fact}

For connected $G$ with $|V(G)|\ge2$ we have $\delta(G)\ge1$, so the statement
is equivalent to the integral inequality
\begin{equation}
 2\,r(G)\le t(G)\,\delta(G).\label{eq:integral}
\end{equation}
(The one-vertex graph makes $\delta=0$; \eqref{eq:integral} then reads $0\le0$
and is trivial, so we always assume $|V(G)|\ge2$.)

The bound is tight: the path $P_{2r}$ has $\delta=1$ and $t=2r$.  Hence the
case $\delta(G)=1$ cannot be improved, and it is the whole difficulty:

\begin{lemma}[Leaf lemma]\label{lem:leaf}
Every finite connected graph with a leaf contains an induced tree on at least
$2\,r(G)$ vertices.
\end{lemma}

Indeed, if $\delta(G)\ge2$ then $t(G)\ge r(G)+1$ by
\cref{lem:geod} below, whence
$t(G)\delta(G)\ge 2(r+1)>2r$; if $\delta(G)=1$ the leaf lemma gives
$t(G)\delta(G)=t(G)\ge 2r$.  This proves
\cref{fact:84} modulo \cref{lem:leaf}, which is proved in
\cref{sec:leaf}.

The proof of the leaf lemma has two branches, according to whether deleting
the leaf lowers the radius.  One branch is governed by the following
classical notion.

\begin{definition}\label{def:vrd}
A connected nontrivial graph $G$ is \emph{vertex-radius-decreasing}
(\emph{vrd}) if $r(G-v)<r(G)$ for every non-cut vertex $v$ of $G$.  By
\cref{lem:drop1} the inequality is automatically an equality:
$r(G-v)=r(G)-1$.
\end{definition}

\begin{definition}\label{def:SlB}
For a connected graph $B$ and $\ell\ge1$, let $S_\ell(B)$ be the graph
obtained from $B$ by attaching to \emph{every} vertex of $B$ a pendant path
of length $\ell$ (a path on $\ell$ new vertices, whose last vertex is a leaf
of $S_\ell(B)$), using that vertex as the first endpoint.
\end{definition}

\begin{theorem}[Structure of vrd graphs with a cut vertex]\label{thm:vrd}
If $G$ is vrd, has a cut vertex, and $r(G)\ge2$, then
$G=S_\ell(B)$ for some $\ell\ge1$ and some connected block $B$
(of order at least two), every vertex of which has degree
$\ell+1$ in $G$; moreover
\[
        r(G)=r(B)+\ell .
\]
Consequently $t(G)\ge 2\,r(G)$
(\cref{cor:vrd-tree}).
\end{theorem}

Classically this is due to Gliviak~\cite{Gliviak1976} (see also the survey
\cite{Gliviak2000}, Theorem~4, and Swart~\cite[Theorem~2.2.10]{Swart1996}),
with ingredients credited to Fajtlowicz~\cite{Fajtlowicz1988}.  Here we
prove \cref{thm:vrd} from scratch in \cref{sec:vrd}.
The non-cut-vertex argument uses paths directly, without a block tree.
The proof of the leaf lemma also uses a \emph{rooted} form of Chung's
induced-path construction (\cref{lem:rooted-chung}), proved in full in
\cref{sec:chung}, including two details that are usually left implicit.

The manuscript is organized as follows.  \cref{sec:elem} collects the
elementary facts, all proved.  \cref{sec:chung} establishes the rooted Chung
lemma, self-contained.  \cref{sec:vrd} develops \cref{thm:vrd} and its
corollary from first principles.  \cref{sec:leaf} proves the leaf lemma from
the two preceding sections, and \cref{sec:main} assembles \cref{fact:84}.
No result of~\cite{ESS1986,Fajtlowicz1988,Gliviak1976,Gliviak2000,Swart1996}
is used in any proof; those works are cited only as historical credits.
The companion Lean development is summarized in \cref{sec:lean}, with a
statement-by-statement correspondence in \cref{app:map}.

\section{Elementary facts}\label{sec:elem}

Throughout, $G$ is a finite connected graph with $n$ vertices; $v$ is a
vertex with $G-v$ connected and nonempty when a radius $r(G-v)$ is mentioned.

\begin{lemma}[Geodesics; the bound $t\ge r+1$]\label{lem:geod}
Every geodesic path induces a path subgraph (no chord joins two of its
vertices non-consecutively).  Consequently
$\rad(G)\le\diam(G)$ and $t(G)\ge\diam(G)+1\ge r(G)+1$.
\end{lemma}

\begin{proof}
A chord would shorten the geodesic.  Vertices at distance $\diam(G)$ are
joined by a geodesic, which is induced, of order $\diam(G)+1$; and
$\rad(G)\le\ecc(x)\le\diam(G)$ for every $x$.
\end{proof}

\begin{lemma}[Radius drops by at most one]\label{lem:drop1}
If $G-v$ is connected and nonempty then $r(G)\le r(G-v)+1$.
\end{lemma}

\begin{proof}
Let $c'$ be a centre of $G-v$ and $\rho=r(G-v)$.  For $x\ne v$,
$d_G(c',x)\le d_{G-v}(c',x)\le\rho$.  The vertex $v$ has a neighbour $w$ in
$G-v$, so $d_G(c',v)\le1+\rho$.  Hence $\ecc_G(c')\le\rho+1$.
\end{proof}

\begin{lemma}[Leaves]\label{lem:leafbasics}
Let $z$ be a leaf of $G$ with $|V(G)|\ge3$.
\begin{enumerate}[label=(\roman*)]
\item Every path with both ends $\ne z$ avoids $z$; hence distances between
      vertices of $G-z$ agree in $G$ and $G-z$.
\item $r(G-z)\le r(G)$: a centre $c$ of $G$ is not a leaf, and
      $\ecc_{G-z}(c)=\ecc_G(c)=r(G)$ by (i).
\item $z$ is not central: $\ecc(z)=1+\ecc(p(z))\ge 1+r(G)>r(G)$.
\item $p(z)$ is a cut vertex (for $|V(G)|\ge3$).
\end{enumerate}
\end{lemma}

\begin{proof}
(i) A path through $z$ must enter and leave $z$ along distinct edges, but $z$
has degree one.  (ii) $z\ne c$ by (iii).  (iii) Clear.  (iv) $z$ is isolated
in $G-p(z)$, while $G-p(z)$ has other vertices.
\end{proof}

\begin{lemma}[Unique eccentric points, I: the drop]\label{lem:uepdrop}
Call $x$ the \emph{unique eccentric point} (\emph{UEP}) of $c$ if
$d(c,x)=\ecc(c)$ and $d(c,y)<\ecc(c)$ for all $y\ne x$.  If $c$ is central
and has UEP $v$, then $G-v$ is connected and $r(G-v)=r(G)-1$.
\end{lemma}

\begin{proof}
Write $r=r(G)$.  For $y\ne v$, a $c$--$y$ path through $v$ has length at
least $d(c,v)+1\ge r+1>r-1\ge d(c,y)$; so every $c$--$y$ geodesic avoids $v$,
and $d_{G-v}(c,y)=d_G(c,y)\le r-1$.  Thus every vertex of $G-v$ reaches $c$:
$G-v$ is connected, and $\ecc_{G-v}(c)\le r-1$.  The reverse inequality is
\cref{lem:drop1}.
\end{proof}

\begin{lemma}[Unique eccentric points, II: the converse]\label{lem:dropuep}
If $v$ is a non-cut vertex and $r(G-v)=r(G)-1$, then some central vertex $c$
of $G$ has UEP $v$.
\end{lemma}

\begin{proof}
Let $c$ be a centre of $G-v$; then $d_G(c,x)\le d_{G-v}(c,x)\le r-1$ for
$x\ne v$, while $\ecc_G(c)\ge r(G)=r$ forces $d_G(c,v)\ge r$.  Also $v$ has a
neighbour $w$ in $G-v$, giving $d_G(c,v)\le1+d_G(c,w)\le r$.  Hence
$d_G(c,v)=r=\ecc_G(c)$: $c$ is central and $v$ is its UEP.
\end{proof}

\begin{lemma}[A vertex has at most one UEP]\label{lem:uepunique}
If $x\ne x'$ are both UEPs of $c$, then $d(c,x')<\ecc(c)$ (from $x$ being the
UEP) contradicts $d(c,x')=\ecc(c)$ (from $x'$ being the UEP).
\end{lemma}

\begin{lemma}[Every connected graph has at least two non-cut vertices]
\label{lem:twoncv}
The first and last vertices of a longest path are non-cut.
\end{lemma}

\begin{proof}
Let $v_1\ldots v_k$ be a longest path and $v=v_1$.  Every neighbour of $v$
lies on it, else the path could be extended.  For $x\ne v$, follow a
geodesic $x=x_0,\dots,x_m=v$; the predecessor $x_{m-1}$ is a neighbour of
$v$, hence equals some $v_i$ with $i\ge2$.  Then
$x_0,\dots,x_{m-1},v_i,v_{i-1},\dots,v_2$ is a walk in $G-v$ from $x$ to
$v_2$.  So $G-v$ is connected.
\end{proof}

\begin{lemma}\label{lem:small}
If $r(G)\ge2$ then $|V(G)|\ge4$.
\end{lemma}

\begin{proof}
$r\ge2$ means no universal vertex; graphs on at most three vertices
all have a universal vertex (on three vertices, both $P_3$ and $K_3$ do).
\end{proof}

\begin{lemma}[Counting; heredity]\label{lem:count}
\begin{enumerate}[label=(\roman*)]
\item For $s,\ell\ge1$:\ $(s+1)(\ell+1)-2(s+\ell)=(s-1)(\ell-1)\ge0$.
\item $t(G-v)\le t(G)$.
\end{enumerate}
\end{lemma}

\begin{lemma}[Auxiliary lemma]\label{lem:aux}
Let $D$ be a finite simple graph with minimum degree at least one, and let
$a\in V(D)$.  If every vertex $v\ne a$ has a neighbour of degree one, then
$a$ also has a neighbour of degree one.
\end{lemma}

\begin{proof}
Suppose $a$ has no neighbour of degree one, and let $x\in N(a)$; then
$\deg(x)\ge2$.  Since $x\ne a$, $x$ has a neighbour $y$ of degree one, so
$N(y)=\{x\}$.  If $y\ne a$, applying the hypothesis to $y$ forces $\deg(x)=1$,
a contradiction; hence $y=a$ and $\deg(a)=1$.  Pick
$z\in N(x)\setminus\{a\}$, and let $w$ be a neighbour of $z$ of degree one
($z\ne a$).  Then $w\ne a$ (as $N(a)=\{x\}\not\ni z$) and $w\ne x$ (as
$\deg(x)\ge2$), so $N(w)=\{z\}$; applying the hypothesis to $w$ forces
$\deg(z)=1$, contradicting $z$ being adjacent to the two distinct vertices
$x$ and $w$.
\end{proof}

\section{The rooted Chung lemma}\label{sec:chung}

\begin{lemma}[Rooted Chung lemma]\label{lem:rooted-chung}
Let $F$ be connected with $r(F)=r\ge2$, and let $a$ be a non-cut vertex with
$r(F-a)=r-1$.  Then $F$ contains an induced path on at least $2r-1$ vertices
that contains $a$.
\end{lemma}

\begin{proof}
Let $v_0$ be a centre of $F-a$, so $\ecc_{F-a}(v_0)=r-1$ and $v_0\ne a$.
Every $x\ne a$ satisfies $d_F(v_0,x)\le d_{F-a}(v_0,x)\le r-1$.

\emph{Claim 0: $d_F(v_0,a)=r$.}
Since $\ecc_F(v_0)\ge r(F)=r$ while all other vertices are at distance at
most $r-1$, we get $d_F(v_0,a)\ge r$.  Conversely $a$ has a neighbour
$w\ne a$ (as $|V(F)|\ge2$ and $F$ is connected), and
$d_F(v_0,a)\le1+d_F(v_0,w)\le r$.  (This uses that $a$ is non-isolated; the
inequality $\le r$ is the first of two details that are easy to omit.)

Fix a geodesic $v_0,v_1,\dots,v_r=a$; it is induced
(\cref{lem:geod}) and $d(v_0,v_j)=j$.

Since $\ecc_F(v_2)\ge r$, there is a vertex $w$ with $d_F(v_2,w)\ge r$; then
$w\ne a$ because $d(v_2,a)=r-2<r$.  Now
\[
 d_F(v_0,w)\ \ge\ d(v_2,w)-d(v_0,v_2)\ \ge\ r-2,\qquad
 d_F(v_0,w)\ \le\ r-1 ,
\]
the latter because $w\ne a$.  Let $P=p_0,\dots,p_m$ be a $v_0$--$w$ geodesic
($p_0=v_0$, $p_m=w$, $m=d(v_0,w)\in\{r-2,r-1\}$); it is induced.

\emph{Chord exclusion.}
(1) For $j\ge2$, no $p_i$ is adjacent to $v_j$: otherwise
\[
 m=d(v_0,w)\ \ge\ (j-1)+\bigl(d(v_j,w)-1\bigr)\ \ge\ (j-1)+\bigl(r-j+2-1\bigr)=r,
\]
using $d(v_j,w)\ge d(v_2,w)-d(v_2,v_j)\ge r-(j-2)$, contradicting $m\le r-1$.
(2) $v_1$ is not adjacent to any $p_i$ with $i\ge2$: otherwise, following
$v_2,v_1,p_i$, then $P$ to $w$, gives
$d(v_2,w)\le 1+1+(m-i)\le m\le r-1<r$, a contradiction.
(3) The single possible chord is $v_1p_1$.
Moreover, \emph{if $v_1p_1$ is an edge then $m=r-1$}: the same detour as in
(2) with $i=1$ gives $r\le d(v_2,w)\le m+1$, i.e.\ $m\ge r-1$.  (This is the
second implicit detail: it is what makes the deletion below lose only one
vertex.)

\emph{Conclusion.}  Concatenate
$Q=v_r,v_{r-1},\dots,v_0,p_1,\dots,p_m$.  By (1) and (2) the only possible
chord of $Q$ is $v_1p_1$, and $a=v_r\in Q$; the order is $(r+1)+m$.
If $v_1p_1$ is not an edge, $Q$ is an induced path of order at least
$(r+1)+(r-2)=2r-1$.  If $v_1p_1$ is an edge, then $m=r-1$, and deleting
$v_0$ from $Q$ leaves $v_r,\dots,v_1,p_1,\dots,p_m$, an induced path of
order $r+(r-1)=2r-1$ (now $v_1p_1$ is a path edge and, by (2), $v_1$ has no
other chord to $P$).  In both cases the path contains $a$, which is never
deleted.
\end{proof}

\begin{remark}
The classical (unrooted) Erd\H{o}s--Saks--S\'os theorem
\cite{ESS1986,ChungCredit} is \emph{not} used in this paper: the case
$\delta\ge2$ of the conjecture needs only \cref{lem:geod}, and in the
vrd branch the core block is handled with the weaker count of
\cref{lem:count}(i).
\end{remark}

\section{The structure theorem for vrd graphs}\label{sec:vrd}

In this section $G$ is vrd (\cref{def:vrd}) with a cut vertex and
$r:=r(G)\ge2$.  Write $\NC(G)$ for the set of non-cut vertices.

\begin{proposition}[Non-cut vertices are leaves]\label{prop:ncvleaf}
Every $v\in\NC(G)$ is a leaf.
\end{proposition}

\begin{proof}
Suppose $v\in\NC(G)$ with $\deg(v)\ge2$.  Let
\[
 S=\{x:\ x\text{ is reachable from }v\text{ by a path whose interior and
 endpoints all lie in }\NC(G)\},
\]
the component of $v$ in $G[\NC(G)]$ ($S$ contains $v$ and is connected by
construction).  We claim that $S$ is closed under neighbours.  Let $w\in S$.
If $\deg(w)\ge2$, then by \cref{lem:dropuep,lem:uepdrop}
(with $w\in\NC$ and $r(G-w)=r-1$) $w$ is the UEP of a central vertex;
then every neighbour $x$ of $w$ satisfies $d(c,x)=r-1$, and any
$c$--$q$ path through $x$ has length at least $(r-1)+1=r>d(c,q)$, so
geodesics from $c$ avoid $x$; $w$ has a second neighbour $y\ne x$,
which reaches $c$ avoiding $x$, so $G-x$ is connected:
$x\in\NC(G)$.  Being adjacent to $w\in S$, we get $x\in S$.
If $\deg(w)=1$, then $w\ne v$ (as $\deg(v)\ge2$), and the $\NC$-path
from $v$ to $w$ has a last edge into $w$; the unique neighbour of $w$
is the penultimate vertex of that path, which lies in $S$.
So no edge leaves $S$.  But $S$ is nonempty ($v\in S$), proper (it contains
only non-cut vertices, while $G$ has a cut vertex), connected, and $G$ is
connected --- a contradiction.
\end{proof}

\begin{corollary}\label{cor:ncvleaf}
The graph $G$ has a leaf: $\NC(G)$ consists of leaves, and $|\NC(G)|\ge2$ by
\cref{lem:twoncv}.
\end{corollary}

\begin{proof}
Immediate from \cref{prop:ncvleaf} and \cref{lem:twoncv}.
\end{proof}

\begin{proposition}[Leaf stripping]\label{prop:strip}
Let $L$ be the set of leaves of $G$ and $G_1=G-L$.  Then $G_1$ is connected
and $r(G_1)=r-1$.  Moreover every vertex of $G_1$ has degree at least two in
$G_1$.
\end{proposition}

\begin{proof}
$G_1$ is connected (deleted vertices are leaves).  Lower bound: a centre $c$
of $G_1$ satisfies $d_G(c,x)=d_{G_1}(c,x)\le r(G_1)$ for non-leaves $x$
(paths avoid leaves, \cref{lem:leafbasics}(i)) and
$d_G(c,z)\le1+r(G_1)$ for leaves $z$; hence
$r(G)\le r(G_1)+1$, i.e.\ $r(G_1)\ge r-1$.  Upper bound: $r(G-u)=r-1$ for a
leaf $u$ (\cref{def:vrd} and \cref{lem:drop1}), and deleting
the remaining leaves one by one never increases the radius
(\cref{lem:leafbasics}(ii)).
\end{proof}

\begin{proposition}[The penultimate layer]\label{prop:penult}
Let $L$, $G_1$ be as above and let $w$ be a non-cut vertex of $G_1$.  Then
$w$ is a cut vertex of $G$ and has at least one leaf neighbour.
\end{proposition}

\begin{proof}
$w\notin L$ (leaves of $G$ are isolated in $G_1$, and $|V(G_1)|\ge2$).  If
$w$ were a non-cut vertex of $G$, \cref{prop:ncvleaf} would make
$w$ a leaf, a contradiction; so $w$ is a cut vertex of $G$.  Now $G-w$ is
disconnected, while $G_1-w$ is connected; so $G_1-w$ lies inside one
component of $G-w$, and every other component consists of leaves only.  A
component of $G-w$ consisting of leaves is a single vertex adjacent to $w$
(two leaves cannot be adjacent: each would have its neighbour outside the
component).  Such components exist because $G-w$ is disconnected, so $w$ has
a leaf neighbour.
\end{proof}

\begin{proposition}[Two leaves do not share a parent]\label{prop:oneleaf}
A vertex of $G$ has at most one leaf neighbour.
\end{proposition}

\begin{proof}
Suppose leaves $u,u'$ share the parent $b$.  By
\cref{lem:dropuep} (applied to the non-cut vertex $u$) some central $c$
has UEP $u$, so $d(c,b)=r-1$ and $c\ne b$ (else $r=d(b,u)=1$).  Every
$c$--$u'$ path ends with $b\,u'$, so $d(c,u')=d(c,b)+1=r$, contradicting the
uniqueness in $\mathrm{UEP}(c,u)$ (as $u'\ne u$).
\end{proof}

\begin{proposition}[Stripping preserves vrd]\label{prop:stripsvrd}
$G_1=G-L$ is vrd: for every non-cut vertex $w$ of $G_1$,
$r(G_1-w)=r(G_1)-1=r-2$.
\end{proposition}

\begin{proof}
Fix such $w$.  By \cref{prop:penult}, $w$ has a leaf neighbour $u$;
by \cref{prop:oneleaf} it is the unique one.  Let $c$ be central
in $G$ with $\mathrm{UEP}(c,u)$ (\cref{lem:dropuep} applied to $u\in\NC(G)$):
$d(c,u)=r$, all others $\le r-1$, so $d(c,w)=r-1$; also $c\notin L$
(\cref{lem:leafbasics}(iii)), so $c\in V(G_1)$ and
$\ecc_{G_1}(c)=r-1=r(G_1)$.

\emph{Claim: $d(c,z)\le r-2$ for all $z\in V(G_1)-w$.}  Suppose
$d(c,z)=r-1$ (distances agree with $G$ by \cref{lem:leafbasics}(i)).
Then $z\notin L$ and $z\ne w$, so $z$ is a cut vertex of $G$ (otherwise
\cref{prop:ncvleaf} puts $z$ in $L$).  Let $C_0\ni c$ be the
component of $G-z$ and let $C$ be another component, nonempty.  For
$y\in C$, every $c$--$y$ path passes through $z$, so
$d(c,y)\ge d(c,z)+1=r$; since $\ecc(c)=r$, we get $d(c,y)=r$, hence $y=u$ by
$\mathrm{UEP}(c,u)$.  So $C=\{u\}$.  But $w$ is adjacent to $u$ and
$w\ne z$, so $w\in C$, i.e.\ $w=u$ --- absurd (as $u$ is a leaf and $w$ has
degree at least two in $G_1$).  This proves the claim.

Hence $\ecc_{G_1-w}(c)\le r-2$, so $r(G_1-w)\le r-2$; the reverse bound is
\cref{lem:drop1} applied in $G_1$.
\end{proof}

\begin{proposition}[Decreasing blocks]\label{prop:block}
Let $B$ be a connected graph with no cut vertex such that
$r(B-v)=r(B)-1$ for every $v$ (e.g.\ the terminal graph of the stripping
chain, by \cref{prop:stripsvrd}).  Then $B$ is self-centered and
admits a fixed-point-free involution $f:V(B)\to V(B)$ with
$\mathrm{UEP}(f(v),v)$ for every $v$; in particular each vertex has a unique
eccentric vertex.
\end{proposition}

\begin{proof}
For each $v$, \cref{lem:dropuep} gives a central $c_v$ with
$\mathrm{UEP}(c_v,v)$; by \cref{lem:uepunique} the map
$f:v\mapsto c_v$ is injective.  Hence $|V|\le|C(B)|\le|V|$, so
$C(B)=V(B)$: $B$ is self-centered, and $f$ is a bijection with $f(v)\ne v$
(as $d(v,f(v))=r(B)\ge1$).  If $f(f(v))\ne v$, then $f(f(v))$ and $v$ are
both at distance $r(B)$ from $f(v)$ (the former because
$\mathrm{UEP}(f^2(v),f(v))$ says $d(f^2(v),f(v))=r(B)$), contradicting
uniqueness of the eccentric vertex of $f(v)$; so $f^2=\mathrm{id}$.  Finally
$\mathrm{UEP}(v,f(v))$ is the statement
$\mathrm{UEP}(f(f(v)),f(v))$ read at $v$.
\end{proof}

\begin{theorem}[Structure, self-contained version of \cref{thm:vrd}]
\label{thm:vrdproof}
$G=S_k(B)$ for the terminal block $B$ of the stripping chain and
$k=r-r(B)\ge1$; in particular $r(G)=r(B)+k$.
\end{theorem}

\begin{proof}
Form the stripping chain $G=G_0\supset G_1\supset\cdots$,
$G_{i+1}=G_i-L(G_i)$ ($L(\cdot)$ the leaf set).  By
\cref{prop:strip,prop:stripsvrd}, every $G_i$ with a
cut vertex is vrd and $r(G_i)=r-i$; since $|V(G_i)|$ strictly decreases,
the chain terminates at the first $k$ with $G_k=:B$ having no cut vertex,
and $r(B)=r-k\ge1$, so $|V(B)|\ge2$.

We show by downward induction on $i$ that $G_i=S_{k-i}(B)$.
For $i=k$ this holds with $S_0(B):=B$.  Assume
$G_{i+1}=S_m(B)$ with $m=k-i-1\ge0$, and let $P$ be the set of parents of
$L(G_i)$; each leaf of $G_i$ has its parent in $G_{i+1}$ (a leaf whose
parent were also deleted would force $G_i=K_2$, whose radius $1$ contradicts
$r(G_i)=r-i\ge r-k+1\ge2$), and each vertex has at most one leaf
(\cref{prop:oneleaf}).  We must show that $P$ is exactly the set
of the $|B|$ \emph{tips} of $G_{i+1}=S_m(B)$ (the vertices at depth $m$,
one on each pendant path).

\emph{(a) Every tip is in $P$.}  Since $r(G_i)=r(B)+m+1$, every $z\in B$
satisfies $\ecc_{G_i}(z)\ge r(B)+m+1$.  Within $G_{i+1}=S_m(B)$, a vertex at
depth $j$ on the branch over $b'$ is at distance at most
$d(z,b')+j\le r(B)+m$ from $z$ (as $d(z,b')\le r(B)$, $B$ being
self-centered).  So some $p\in P$ has $d(z,p)\ge r(B)+m$; also
$d(z,p)\le r(B)+\mathrm{depth}(p)\le r(B)+m$.  Hence $d(z,p)=r(B)+m$ and the
witness is unique: $B$ has a unique eccentric vertex $f(z)$
(\cref{prop:block}), and the distance bound shows
$p$ lies on $f(z)$'s branch at depth $m$, i.e.\ $p$ is $f(z)$'s tip.  As
$f$ is surjective, every tip is in $P$.

\emph{(b) No other vertex is in $P$.}  Suppose $p\in P$ has depth $j<m$ on
the branch over $b\in B$, with leaf $u$ adjacent to $p$.  The vertex $u$ is
a non-cut vertex of $G_i$, so by \cref{lem:dropuep} some centre $c$ of
$G_i$ has $\mathrm{UEP}(c,u)$: $d(c,u)=r(B)+m+1$, hence $d(c,p)=r(B)+m$.
Now $c$ is not a leaf of $G_i$ (\cref{lem:leafbasics}(iii)).  If
$c\in B$, then $d(c,p)\le r(B)+j\le r(B)+m-1$, too small.  If $c$ lies on
$p$'s own branch at depth $j'$, then
$d(c,p)=|j'-j|\le m$, too small.  If $c$ lies on the branch over
$b''\ne b$ at depth $j''$, let $w$ be the tip
of $p$'s branch (depth $m$; it exists and lies in $G_{i+1}$).  Then
\[
 d(c,w)=j''+d(b'',b)+m=d(c,p)+(m-j)\ \ge\ r(B)+m+1=\ecc(c),
\]
so $d(c,w)=\ecc(c)$, and uniqueness in $\mathrm{UEP}(c,u)$ forces $w=u$;
but $w\in V(G_{i+1})$ while $u\notin V(G_{i+1})$, a contradiction.

This completes the induction: $G=G_0=S_k(B)$, and
$r(G)=r(B)+k$ because the stripping chain drops the radius by exactly one
per step (\cref{prop:strip,prop:stripsvrd}).
Moreover every vertex of $B$ has degree $\ell+1$ in $G$ for the appropriate
layer, as asserted in \cref{thm:vrd}.
\end{proof}

\begin{corollary}[$t\ge 2r$ for vrd graphs with a cut vertex]
\label{cor:vrd-tree}
If $G$ is vrd with a cut vertex, then $t(G)\ge2\,r(G)$.
\end{corollary}

\begin{proof}
Write $G=S_k(B)$ as in \cref{thm:vrdproof}, $s=r(B)\ge1$.
Take a geodesic $Q$ of $B$ (order at least $s+1$,
\cref{lem:geod}).  The union of $V(Q)$ with, for each $q\in Q$, the
$\ell$-vertex pendant path attached to $q$, induces a tree: pendant paths
are internally disjoint, there are no edges between distinct pendant paths,
and $Q$ is induced in $B$ (\cref{lem:geod}); its order is at least
\[
 (s+1)(k+1)\ \ge\ 2(s+k)\ =\ 2\,r(G)
\]
by \cref{lem:count}(i).
\end{proof}

\begin{remark}
The case $B=K_2$ needs no special treatment: $K_2$ satisfies
\cref{prop:block} with $f$ the transposition, and equality holds
throughout \cref{cor:vrd-tree} ($G=P_{2k+2}$, $t=2r$).
\end{remark}

\begin{remark}[Historical credit]
Theorems of this shape are classical: Gliviak~\cite{Gliviak1976} (see
\cite{Gliviak2000}, Theorem~4, and Swart~\cite[Theorem~2.2.10]{Swart1996})
describes radially critical graphs as a ``decreasing'' core with equilength
endpaths, and the fact that non-cut vertices of vrd graphs with a cut vertex
are leaves is attributed to Fajtlowicz~\cite{Fajtlowicz1988}.  The proofs
above were reconstructed from first principles and are self-contained; we
did not consult the original arguments.
\end{remark}

\section{The leaf lemma}\label{sec:leaf}

\begin{proof}[Proof of \cref{lem:leaf}]
Suppose the lemma fails and let $G$ be a counterexample of minimum order:
$G$ connected with a leaf $u$ and $t(G)<2r(G)$, while every smaller connected
graph with a leaf satisfies the lemma.  Let $a=p(u)$ and $r=r(G)$.

If $r=1$, the edge $ua$ induces a tree on $2=2r$ vertices, a contradiction.
So $r\ge2$; by \cref{lem:small} $|V(G)|\ge4$, and $a$ is a cut vertex
(\cref{lem:leafbasics}(iv)).

\emph{Claim 1.} \emph{For every non-cut vertex $v\ne u$ of $G$:
$r(G-v)=r-1$.}
Indeed $v\ne a$ (as $a$ is a cut vertex and $v$ is not), so $G-v$ is
connected and still has the leaf $u$.  If $r(G-v)\ge r$, minimality
gives $t(G-v)\ge2r(G-v)\ge2r$, while $t(G-v)\le t(G)$
(\cref{lem:count}(ii)) --- a contradiction.  So $r(G-v)<r$, and
\cref{lem:drop1} gives $r(G-v)=r-1$.

\emph{Setup.} Let $H=G-u$.  By \cref{lem:leafbasics}, distances among
$V(H)$ agree in $G$ and $H$, and $u$ is not central; hence centres of
$G$ survive in $H$ and $r(H)\le r$.  If $r(H)<r$, then $u$ is a
non-cut vertex and \cref{lem:drop1} gives $r(H)=r-1$.  So
$r(H)\in\{r-1,\,r\}$.

\medskip
\noindent\textbf{Case A} ($r(H)=r-1$).
Then every non-cut vertex of $G$ deletes the radius to $r-1$: $u$ by the
case assumption, all others by Claim 1.  So $G$ is vrd with the cut vertex
$a$, and \cref{cor:vrd-tree} gives $t(G)\ge2r$ --- a contradiction.
(Minimality is not needed in this case.)

\medskip
\noindent\textbf{Case B} ($r(H)=r$).

\emph{Claim 2.} \emph{$H$ has no leaf; in particular $\delta(H)\ge2$.}
Otherwise $H$ is a smaller connected graph with a leaf and
$r(H)=r$, so minimality gives $t(H)\ge2r(H)=2r$; but
$t(G)\ge t(H)$ (\cref{lem:count}(ii)), contradiction.

\emph{Claim 3.} \emph{For every non-cut vertex $v\ne a$ of $H$:
$r(H-v)=r-1$.}
$G-v$ is obtained from the connected graph $H-v$ by re-attaching the
leaf $u$ at $a$, hence is connected; so $v$ is a non-cut vertex of
$G$, and Claim 1 gives $r(G-v)=r-1$.  The graph $G-v$ has at least
three vertices, so its leaf $u$ is not central
(\cref{lem:leafbasics}(iii)); hence centres of $G-v$ survive in
$H-v$ and $r(H-v)\le r-1$.  The reverse inequality is
\cref{lem:drop1} (as $r(H)=r$ and $v$ is a non-cut vertex of
$H$).

\emph{Claim 4.} \emph{$H$ has no cut vertex.}
Claims 2 and 3 and \cref{lem:dropuep} imply that every non-cut
vertex $v\ne a$ is a central vertex's UEP and has degree at least two.
The neighbour argument in \cref{prop:ncvleaf} therefore
shows that every neighbour of $v$ is also a non-cut vertex.
Choose a non-cut vertex $v\ne a$; two distinct non-cut vertices exist
by the longest-path argument.
If $a$ were a cut vertex, the non-cut property would propagate along
any path starting at $v$: a non-cut vertex can never equal $a$, so
the neighbour implication applies at every step. A path from $v$
to $a$ would give a contradiction. Thus $a$ is a non-cut vertex.
For any $w\ne a$, choose a $v$--$w$ path in the connected graph $H-a$.
Propagating the same implication along this path shows that $w$
is a non-cut vertex. Hence every vertex of $H$ is a non-cut vertex.

By Claim 4 every vertex of $H$ is a non-cut vertex.  For each $v\ne a$,
Claim 3 and \cref{lem:dropuep} provide a central $c_v$ with
$\mathrm{UEP}(c_v,v)$; distinct $v$ give distinct $c_v$
(\cref{lem:uepunique}).  So with $n=|V(H)|$, at least $n-1$ vertices of
$H$ are central.

\emph{Claim 5.} \emph{$H$ is self-centered.}
If some $b$ were non-central, then every $y\ne b$ is central, whence
$d(y,b)\le r$, so $d(b,y)\le r$ and $\ecc(b)\le r$; but always
$\ecc(b)\ge r(H)=r$.  Hence $\ecc(b)=r$ and $b$ is central, a
contradiction.

Form the auxiliary graph $D$ on $V(H)$ by putting $xy\in E(D)$ iff
$d_H(x,y)=r$.  Since $H$ is self-centered and finite, every vertex has a
vertex at distance exactly $r$: $\delta(D)\ge1$.  For each $v\ne a$, $c_v$
is adjacent to $v$ in $D$ and $\deg_D(c_v)=1$ (any other $w$ with
$d(c_v,w)=r$ would violate $\mathrm{UEP}(c_v,v)$).  So every $v\ne a$ has a
degree-one neighbour, and \cref{lem:aux} provides $c$ adjacent to $a$
in $D$ with $\deg_D(c)=1$.  Thus $d_H(c,a)=r$ and $a$ is the unique vertex
at distance $r$ from $c$; by Claim 5, $c$ is central, so
$\mathrm{UEP}(c,a)$.  \cref{lem:uepdrop} gives $r(H-a)=r-1$, and $a$
is a non-cut vertex of $H$ (Claim 4).  \cref{lem:rooted-chung} yields
an induced path $P$ of $H$ with $|P|\ge2r-1$ and $a\in P$.

Finally, $u$ is adjacent in $G$ only to $a\in P$, so $G[V(P)\cup\{u\}]$ is
an induced tree on at least $2r$ vertices --- the final contradiction.
\end{proof}

\section{Graffiti.pc Conjecture 84}\label{sec:main}

\begin{theorem}\label{thm:main}
Let $G$ be a finite connected graph with at least two vertices.  Then
\[
        2\,r(G)\ \le\ t(G)\,\delta(G),
\]
equivalently $t(G)\ge 2r(G)/\delta(G)$.
\end{theorem}

\begin{proof}
A connected nontrivial graph has $r(G)\ge1$ and $\delta(G)\ge1$.
If $\delta(G)\ge2$, \cref{lem:geod} gives $t(G)\ge r(G)+1$, so
$t(G)\delta(G)\ge 2(r(G)+1)>2r(G)$.
If $\delta(G)=1$, $G$ has a leaf, and \cref{lem:leaf} gives
$t(G)\ge2r(G)$, so $t(G)\delta(G)=t(G)\ge2r(G)$.
\end{proof}

\begin{remark}
The leaf lemma is genuinely about induced trees: one cannot strengthen it to
an induced \emph{path} on $2r$ vertices.  In Case B of its proof the extra
vertex is a pendant attachment to a long induced path; in Case A the tree is
assembled from several pendant branches over a geodesic of the core block.
\end{remark}

\section{Supplementary Lean 4 formalization}\label{sec:lean}

An accompanying Lean~4/Mathlib development formalizes the leaf lemma and the
final theorem together with their principal ingredients, under the pinned
toolchain Lean and Mathlib~4.32.2.  For Case A the formal proof deliberately
takes a shorter route than this manuscript: direct induction by peeling
leaves off a maximum induced tree, without the structural classification of
\cref{sec:vrd}.  The continuous-integration build of the repository compiles
every module on a clean runner, scans the sources for proof placeholders and
custom axioms, and prints the axiom dependencies of the final declarations,
which are the three standard Lean axioms only.  The formalization is
supplementary: the mathematical proof in this article is self-contained, and
\cref{app:map} provides a navigation map between the two developments.

\begin{thebibliography}{9}

\bibitem{ESS1986}
P.~Erd\H{o}s, M.~Saks, and V.~T.~S\'os,
\emph{Maximum induced trees in graphs},
J. Combin. Theory Ser. B \textbf{41} (1986), 61--79.
\url{https://doi.org/10.1016/0095-8956(86)90028-6}

\bibitem{ChungCredit}
F.~R.~K.~Chung,
proof of the induced-path bound, Theorem~2.2 of \cite{ESS1986};
\cref{lem:rooted-chung} is a rooted variant of her construction, with the
proof written out in full.

\bibitem{Fajtlowicz1988}
S.~Fajtlowicz,
\emph{A characterization of radius-critical graphs},
J. Graph Theory \textbf{12} (1988), 529--532.
\url{https://doi.org/10.1002/jgt.3190120409}

\bibitem{Gliviak1976}
F.~Gliviak,
\emph{On the structure of radially critical graphs},
in \emph{Graphs, Hypergraphs and Block Systems},
Proc. Sympos. Zielona G\'ora, Poland (1976), 69--73.

\bibitem{Gliviak2000}
F.~Gliviak,
\emph{On radially extremal graphs and digraphs, a survey},
Math. Bohem. \textbf{125} (2000), 215--225.

\bibitem{Swart1996}
C.~S.~Swart,
\emph{Distance Measures in Graphs and Subgraphs},
M.Sc.\ thesis, University of Natal, 1996.

\bibitem{WestGraffiti}
D.~B.~West,
\emph{Some conjectures of Graffiti.pc},
online catalogue, entry 84 (2004).

\end{thebibliography}

\appendix
\section{Formal theorem map}\label{app:map}

This map is provided for navigation of the supplementary source and is not
part of the logical dependency of the manuscript.  The Case~A branch of the
Lean proof uses the shorter leaf-peeling induction mentioned in
\cref{sec:lean}; the remaining rows correspond to statements proved in both
developments.  The principal correspondences are:

\begin{center}
\small
\begin{tabular}{@{}P{5.6cm}P{9.0cm}@{}}
\toprule
Mathematical statement & Lean declaration / file \\
\midrule
Geodesics induce paths; $t\ge r+1$ &
\texttt{two\_radius\_le\_treeNumber\_mul\_minDegree}\newline
\texttt{GraphConjecture84.lean} \\
Radius drops by at most one &
\texttt{radOn\_le\_radOn\_erase\_add\_one}\newline
\texttt{BasicFacts.lean} \\
UEP I: the drop &
\texttt{isUniqueEccentricPoint\_of\_induce\_radius\_drop}\newline
\texttt{Deletion.lean} \\
UEP II: the converse &
\texttt{induce\_radius\_drop\_of\_central\_uep}\newline
\texttt{Deletion.lean} \\
A vertex has at most one UEP &
\texttt{isUniqueEccentricPoint\_unique}\newline
\texttt{BasicFacts.lean} \\
Two non-cut vertices &
\texttt{exists\_two\_deleteConnected}\newline
\texttt{EndBlocks.lean} \\
Unique leaf neighbour &
\texttt{unique\_leaf\_neighbor}\newline
\texttt{EndBlocks.lean} \\
Auxiliary degree-one lemma &
\texttt{degree\_one\_neighbor\_of\_all\_other}\newline
\texttt{RadiusCriticalStructure.lean} \\
Rooted Chung lemma &
\texttt{rooted\_chung}\newline
\texttt{RootedChung.lean} \\
Case B: non-cut closure &
\texttt{deleteConnected\_all\_of\_drop\_except}\newline
\texttt{EndBlocks.lean} \\
Case B: all vertices central; UEP for $a$ &
\texttt{isCentral\_all\_of\_uep\_except},\newline
\texttt{central\_uep\_of\_uep\_except}\newline
\texttt{CaseB.lean} \\
Case B: rooted induced tree &
\texttt{rooted\_tree\_of\_drop\_except}\newline
\texttt{CaseB.lean} \\
Case A branch (peeling route) &
\texttt{vrd\_cut\_tree\_bound}\newline
\texttt{CaseA.lean} \\
Leaf lemma &
\texttt{leafLemma}\newline
\texttt{LeafLemma.lean} \\
Main theorem &
\texttt{graphConjecture84}\newline
\texttt{GraphConjecture84.lean} \\
\bottomrule
\end{tabular}
\end{center}

\end{document}
